% Generated draft: requires mathematical and source review.
\documentclass[11pt]{article}
\usepackage[margin=25mm]{geometry}
\usepackage{fontspec}
\setmainfont{Latin Modern Roman}
\usepackage{amsmath,amssymb,mathtools,unicode-math}
\setmathfont{Latin Modern Math}
\usepackage{enumitem,needspace,xurl,hyperref}
\setlength{\parindent}{0pt}
\setlength{\parskip}{5pt}
\setlength{\emergencystretch}{3em}
\newcommand{\N}{\mathbb{N}}
\newcommand{\Z}{\mathbb{Z}}
\newcommand{\Q}{\mathbb{Q}}
\newcommand{\R}{\mathbb{R}}
\newcommand{\C}{\mathbb{C}}
\newcommand{\F}{\mathbb{F}}
\newcommand{\A}{\mathbb{A}}
\newcommand{\PP}{\mathbb{P}}
\newcommand{\E}{\mathbb{E}}
\newcommand{\eps}{\varepsilon}
\newcommand{\dd}{\,\mathrm{d}}
\begin{document}
\section*{UnsolvedMath ID: 20001134}
\textbf{Source problem code:} catalogue entry
\subsection{Definitions and mathematical statement}
For a finite group $G$, call $P\subseteq G$ product-free when $P\cap PP=\varnothing$, and define $\nu(G)=\max\{|P|:P\subseteq G\text{ product-free}\}$. Repeated factors are allowed in $PP$. V. S\'os's source asks how large such a set can be. Its opening says ``abelian,'' but its subsequent alternating-group example extends beyond abelian groups; the two settings must be distinguished.

For nontrivial finite abelian $G$, the source attributes to Alon--Kleitman the sharp universal existence bound $\nu(G)\geq2|G|/7$. This is a lower bound for the maximum, not for every product-free subset. The trivial group must be excluded from that formulation. The aid reports that Diananda--Yap and Green--Ruzsa give the exact abelian classification, whose case formulas are not supplied here.

For the alternating group $A_d$, Green's question is whether
\[
 \nu(A_d)=o(|A_d|)\qquad(d\to\infty).
\]
The degree $d$ is distinct from the group order $|A_d|$. The supplied aid says Gowers proved this using quasirandom-group mixing, Eberhard improved the bound, and Keevash--Lifshitz--Minzer classified maximum product-free sets for sufficiently large $d$, with a 2024 classification reference.

The aid's finite-parameter analysis concerns the Crane--Kedlaya family. An explicit form is, for $x\in[d]$ and $I\subseteq[d]\setminus\{x\}$ with $|I|=k$,
\[
 F_I^x=\{\pi\in A_d:\pi(x)\in I,\ \pi(I)\cap I=\varnothing\}.
\]
With composition $(\pi\sigma)(x)=\pi(\sigma(x))$, these sets are product-free. For $d\geq6$ and $1\leq k\leq\lfloor d/2\rfloor$, the reported count and consecutive ratio are
\[
 b_d(k)=|F_I^x|=\frac{k(d-k)!(d-k-1)!}{2(d-2k)!},\qquad
 \frac{b_d(k+1)}{b_d(k)}
 =R_d(k)=\frac{k+1}{k}\frac{(d-2k)(d-2k-1)}{(d-k)(d-k-1)},
\]
with the ratio used only where $k+1$ is in range. The aid claims strict decrease of these ratios, no adjacent tie for every $d\geq6$, and unique optimum $k_d=\lceil\rho_d\rceil$, where $\rho_d$ is the unique real root of $R_d(\rho)=1$ in $(0,(d-1)/2)$. It records
\[
 \rho_d=\sqrt{d/2}-5/8+O(d^{-1/2}).
\]
Combined with the cited classification, it reports $\nu(A_d)=b_d(k_d)$ for sufficiently large $d$. Optimization within this family for $d\geq6$ is distinguished from global optimality, whose classification threshold is unspecified. These are supplied literature and research-aid claims requiring independent verification.
\subsection{Short English statement}
Determine maximum product-free subsets, distinguishing finite abelian groups from alternating groups. The supplied aid reports the alternating-group density conjecture solved and a large-degree classification, then optimizes the explicit Crane--Kedlaya family; the family result alone is not a global theorem for every small degree.
\subsection{Sources}
\begin{itemize}
\item \emph{Problems Presented at the Workshop on Recent Trend in Additive Combinatorics}, collected by Ernie Croot and Vsevolod F. Lev, problem 4.3 attributed to V. S\'os, with Alon--Kleitman's bound and Green's conjecture: \url{https://aimath.org/WWN/additivecomb/additivecomb.pdf}.
\item The attached aid names Diananda--Yap, Green--Ruzsa, Gowers, Eberhard, Crane--Kedlaya, and Keevash--Lifshitz--Minzer (classification described as 2024), without complete bibliographic entries.
\item Local registry \texttt{data/unsolvedmath/problems.json}, permanent ID 20001134, exact source code \texttt{catalogue entry}; canonical location \texttt{aim-combinatorics-notes.json notes[258]}. Attached attempt 1, partial result, problem status at run solved, rigor proved, low novelty confidence; machine-generated, not peer reviewed, requiring independent verification before citation.
\end{itemize}
\end{document}
